\documentclass{situationist-preprint}

\title{Moral Intuition as Perception Under Imperfect Conditions}
\msid{MTE-08-a}
\msfamily{MTE-08}
\msdiscipline{Metaethics}
\mstype{Principal manuscript}
\mskeywords{moral intuition; moral perception; moral knowledge; reliability; safety; framing effects; a posteriori}

\begin{document}
\maketitle

\begin{abstract}
How do we know moral truths? We argue that we know them through intuition, understood not as rational insight into self-evident principles but as a faculty of reaction to particular cases, comparable to perception under imperfect conditions. A moral intuition is a seeming that an act or situation has a moral property, formed non-inferentially in response to the presentation of its other features. By phenomenal conservatism, such seemings confer prima facie justification. Given that moral facts are natural facts about idealized reactions, intuitions are actual reactions of the same kind, and their reliability is the degree to which the conditions in which they are formed approach the idealized ones. Where the conditions are good enough, true intuitions are safe, and so yield knowledge; where they are poor, through ignorance, partiality or the framing of the case, intuitions still justify but may fail to yield knowledge, as perception does in bad light. Moral knowledge so acquired is a posteriori, which reconciles intuitionism with the restricted scope of a priori knowledge.
\end{abstract}

\section{The problem}

Moral knowledge seems to be common. We know that torturing a child for amusement is wrong, that one should keep one's promises in ordinary circumstances, and that the slave trade was an injustice. How do we know these things? The question is pressing because the most obvious answers fail. We do not know them by observation in the ordinary sense, since wrongness is not seen as color is seen; we do not know them by deduction from non-moral premises, since MTE-10 argues that an ought can be derived from an is only with a normative bridge premise; and we do not know them by testimony alone, since testimony transmits knowledge but does not generate it (EPI-08).

The intuitionist tradition answers that moral truths are known by intuition. \citet{Sidgwick1907} held that some moral principles are self-evident and laid down conditions under which an apparent self-evidence can be trusted: the proposition must be clear and precise, must survive careful reflection, must be consistent with other such propositions, and must command the agreement of competent judges. \citet{Ross1930} held that the principles of prima facie duty are self-evident in the way that the axioms of mathematics are, known by reflection once one has reached sufficient mental maturity. \citet{Audi2004} and \citet{Huemer2005} have revived the view, Audi with a moderate account of self-evidence on which self-evident propositions can be known on the basis of adequate understanding but need not be obvious, and Huemer with an account of intuitions as intellectual seemings that confer justification by phenomenal conservatism.

A second family of views holds that moral knowledge is, at least in part, perceptual. \citet{McGrath2004} argued that we know moral facts about particular cases in the way we know other facts by perception, \citet{Cullison2010} that there are moral perceptual seemings, \citet{McBrayer2010} that a limited form of moral perception is defensible, and \citet{Audi2013} that moral perception is genuine perception of non-moral properties integrated with a moral response. Against these, \citet{Vayrynen2018} doubts that the contents of perception include moral properties.

The intuitionist answer faces a scope problem and a reliability problem. The scope problem is that EPI-12 argues that a priori knowledge extends only to logic, arithmetic and some conceptual truths, and that intellectual seemings about matters the concepts do not settle can justify but cannot yield knowledge, because they arise in the same way whether or not they are true. If moral intuition is rational insight into non-conceptual principles, it falls under this verdict. The reliability problem is that moral intuitions are affected by factors that seem irrelevant to their truth: the wording of a case, the order in which cases are presented, and the subject's emotional state \citep{Petrinovich1996,SinnottArmstrong2008}. \citet{Haidt2001} argued that moral judgments are typically produced by quick affective responses and rationalized afterwards.

We argue that both problems are solved by a single account of what intuition is. Moral intuition is a faculty of reaction to particular cases, presented in experience, testimony or imagination; its reliability is that of a faculty operating under conditions that are imperfect but often good enough; and the knowledge it yields is a posteriori. Intuition, so understood, is fallible but generally reliable, and comparable to perception under imperfect conditions.

\section{Framework}

\begin{definition}[Moral intuition]\label{d:intuition}
A \emph{moral intuition} is a seeming, in the sense of EPI-19, that a particular act, person or situation has a moral property, formed non-inferentially in response to the presentation of its non-moral features.
\end{definition}

The case may be presented in perception, as when one sees a man beating a dog; in testimony, as when one reads a description of a massacre; or in imagination, as when one considers a hypothetical case. In each, the features of the case are presented, and the moral seeming is a response to them. The definition excludes seemings whose content is a general principle; we return to principles in Section~\ref{s:principles}.

\begin{principle}[Phenomenal conservatism]\label{p:pc}
If it seems to a subject that $p$, she has prima facie justification for believing that $p$.
\end{principle}

\ref{p:pc} is the conclusion of EPI-19. It applies to moral intuitions as to any seemings.

\begin{principle}[Knowledge]\label{p:know}
A subject knows that $p$ if and only if she has a justified belief that $p$ that is safe: in all close cases in which she believes on the same basis, what she believes is true.
\end{principle}

\ref{p:know} is the analysis of EPI-01. For beliefs about particular cases, the close cases are those in which the subject forms a belief by the same faculty, in response to a case presented in a similar way, under similar conditions.

\begin{assumption}[Moral facts as idealized reactions]\label{a:ideal}
Moral facts are natural facts about idealized reactions: an act has a moral property if and only if the reactions to it under idealized conditions, of full information about the case, vivid appreciation of what it means for each person affected, and equal attention to the interests of all of them, would be of the corresponding kind.
\end{assumption}

Assumption~\ref{a:ideal} is the thesis of MTE-01, stated in the form we need. The reactions in question are favoring and disfavoring responses to the features of a case; a moral intuition is a seeming that is, or is grounded in, such a reaction, formed under the conditions that actually obtain.

\begin{definition}[Conditions]\label{d:conditions}
The \emph{conditions} of a moral intuition are the respects in which the circumstances of its formation approach or depart from the idealized conditions: how much the subject knows about the case, how vividly she appreciates its significance for those affected, and how equally she attends to their interests. Conditions are \emph{good enough} for a case if departures of that size do not change the reaction to cases of that kind.
\end{definition}

\begin{assumption}[Continuity]\label{a:cont}
Actual reactions vary continuously with the conditions: the closer the conditions of a reaction to the idealized ones, the more likely it is to agree with the idealized reaction, and for many cases there are conditions short of the ideal that are good enough.
\end{assumption}

Assumption~\ref{a:cont} is the moral analogue of a familiar fact about perception. Color perception in daylight at moderate distance agrees with the colors of things; in dim light or at a great distance it agrees less often; and for many objects, ordinary light is good enough, in that no improvement in the conditions would change the verdict. That the slave trade was wrong is a verdict that no increase in information or impartiality would change; that a particular tax policy is just may be one that it would.

\section{The argument}

\begin{lemma}[Justification]\label{l:just}
Every moral intuition confers prima facie justification for believing its content.
\end{lemma}

\begin{proof}
A moral intuition is a seeming (Definition~\ref{d:intuition}), and by \ref{p:pc} every seeming confers prima facie justification.
\end{proof}

\begin{lemma}[Reliability]\label{l:rel}
A moral intuition formed under conditions good enough for its case agrees with the moral fact, and so do the intuitions formed on the same basis in close cases.
\end{lemma}

\begin{proof}
By Assumption~\ref{a:ideal}, the moral fact about the case is the idealized reaction to it. By Definition~\ref{d:conditions}, under conditions good enough for the case, the actual reaction does not differ from the idealized one, so the intuition agrees with the moral fact. In close cases the subject reacts by the same faculty to cases of the same kind under similar conditions; similar conditions are, for cases of that kind, good enough as well, since goodness of conditions is defined over departures of a size and not over a single case, and so the intuitions formed in those cases agree with the idealized reactions, which are the moral facts.
\end{proof}

\begin{theorem}[Moral knowledge by intuition]\label{t:main}
A true belief formed on the basis of a moral intuition under conditions good enough for its case, in the absence of defeaters, is knowledge.
\end{theorem}

\begin{proof}
The belief is justified by Lemma~\ref{l:just}, since the prima facie justification is undefeated. It is safe by Lemma~\ref{l:rel}, since in all close cases in which the subject believes on the same basis, what she believes agrees with the moral fact. By \ref{p:know}, it is knowledge.
\end{proof}

\begin{corollary}[Fallibility]\label{c:fall}
A belief formed on the basis of a moral intuition under conditions that are not good enough for its case is prima facie justified, but may be false, and even if true may fail to be knowledge.
\end{corollary}

\begin{proof}
It is justified by Lemma~\ref{l:just}. If the conditions are not good enough, departures of that size do change reactions to cases of that kind, so in close cases the subject's intuition may disagree with the idealized reaction, and the belief may be unsafe even if true. By \ref{p:know} it may then fail to be knowledge.
\end{proof}

Theorem~\ref{t:main} and Corollary~\ref{c:fall} together are the comparison with perception. A perceptual belief formed in good light is knowledge if true; one formed in dim light is justified, may be false, and may be unsafe even if true. Moral intuition is \emph{generally reliable} because the conditions in which human beings form their moral intuitions are, for a large class of cases, good enough: the cases of everyday moral life, in which one knows what one is doing and to whom, and in which the interests at stake are visible. It is \emph{fallible} because, for other cases, they are not.

\begin{corollary}[A posteriori]\label{c:apost}
Moral knowledge acquired by intuition is a posteriori.
\end{corollary}

\begin{proof}
The intuition is a response to the presentation of the non-moral features of a case, in perception, testimony or imagination, and the presentation figures in the support of the belief: the subject's justification and the safety of her belief both depend on the features of the case as presented and on the conditions in which she responds to them. In the sense of EPI-12, experience figures as evidence, and the knowledge is not a priori. Where the case is imagined, the features are stipulated, but the subject's knowledge of how cases with those features are to be judged rests on a faculty whose reliability is a contingent fact about the conditions of its exercise, and the safety of the resulting belief depends on that fact.
\end{proof}

\section{Consistency with the scope of the a priori}\label{s:principles}

EPI-12 argues that an intellectual seeming about a matter the concepts do not settle can justify but cannot make knowledge, because such seemings are produced in the same way whether their contents are true or false, and many of them have been false. Corollary~\ref{c:apost} shows that moral intuitions, as Definition~\ref{d:intuition} defines them, are not of that kind. Two differences matter.

First, a moral intuition is a response to the features of a case, and those features are what make the moral fact: by Assumption~\ref{a:ideal}, the moral fact is the idealized reaction to the same features. Had the features been different, the moral fact would have been different, and, under good enough conditions, so would the intuition. The seemings of EPI-12 are not like this, since nothing in their production varies with the truth of their contents.

Second, the reliability of moral intuition is not guaranteed by the concepts. It is a contingent fact about the conditions in which reactions are formed, and it is established, so far as it is, by the kind of evidence that establishes the reliability of perception: agreement among observers in good conditions, stability under improvement of the conditions, and the explanation of errors by identifiable distortions. This is the evidence that MTE-08-b, on calibration, takes up.

What, then, of moral principles? EPI-12 holds that moral principles beyond those that are conceptual are not known a priori; it allows that the rationalist's seemings about them can justify. We agree on both counts. That one ought to keep promises in ordinary circumstances is known, if it is known, by generalization from case intuitions and by the adjustment of cases and principles to each other that \citet{Rawls1951} described and \citet{Daniels1979} developed as wide reflective equilibrium. This is inductive knowledge, a posteriori because its basis is. Ross's view that the principles of prima facie duty are self-evident in the manner of mathematical axioms is, on our account, an overstatement of a truth: the principles seem true on reflection, the seeming justifies, and the principles are known, but they are known through the cases and not by insight into the principles themselves.

\section{The comparison with perception}

The comparison of moral intuition with perception has four points of similarity and one important disanalogy.

\emph{Non-inferential.} A competent observer who sees boys setting a cat on fire does not infer that what they are doing is wrong \citep{Harman1977}; it seems wrong to her, immediately, as the cat seems to be a cat.

\emph{Presentational.} The intuition presents its content as true, with the assertive character that EPI-19 attributes to all seemings, and it can persist when the subject believes it false, as the Müller-Lyer lines continue to look unequal. A subject who has been persuaded that a harmless but disgusting act is permissible may continue to have the intuition that it is wrong.

\emph{Condition-sensitive.} Perception is reliable in some conditions and not in others, and the conditions can be specified: light, distance, the state of the sense organ. Moral intuition is reliable in some conditions and not in others, and the conditions are those of Definition~\ref{d:conditions}.

\emph{Correctable.} Perceptual error is corrected by other perceptions and by theory; moral error is corrected by other intuitions, by reflection on principles, and by improvement of the conditions, as when one learns more about the people affected. The expansion of the moral circle that MTE-19 treats as the measure of moral progress is, on this account, a sustained improvement in one of the conditions, the equality of attention to the interests of those affected.

The disanalogy is that there is no organ of moral intuition. Moral intuitions are not produced by a dedicated sense but by the general capacity to react to what one apprehends; whether they are strictly perceptual, a question on which \citet{McGrath2004} and \citet{Vayrynen2018} disagree, does not matter for the thesis. The thesis says that intuition is \emph{comparable} to perception, and the comparison concerns the epistemology, not the mechanism.

\section{Objections and replies}

\begin{objection}[Framing effects]
Moral intuitions vary with the wording of a case and the order of presentation \citep{Petrinovich1996}. \citet{SinnottArmstrong2008} argues that, since a subject cannot tell from the inside whether her intuition is subject to a framing effect, no moral intuition is justified without inferential confirmation.
\end{objection}

\begin{reply}
Framing effects are departures from good conditions: the wording of a case directs attention to some features and away from others, which is a failure of equal attention. Corollary~\ref{c:fall} predicts that intuitions formed under such conditions may be false and unsafe, and the evidence confirms the prediction for cases of the kinds studied, which are typically unfamiliar, high-stakes dilemmas. It does not show that intuitions in good conditions are unreliable, and the size of the effects varies greatly with the case. Nor does it show that inferential confirmation is needed for justification: \ref{p:pc} gives prima facie justification, which a known framing effect defeats. Perception is subject to illusions that the subject cannot detect from the inside, and we do not conclude that perceptual beliefs need inferential confirmation in good light.
\end{reply}

\begin{objection}[Disagreement]
If intuition were reliable, competent judges would agree. They disagree widely, across cultures and within them.
\end{objection}

\begin{reply}
Agreement is to be expected where conditions are good, and disagreement where they are not; the account predicts disagreement over hard cases and over cases in which information and impartiality are scarce, and agreement over the central cases, which is what we find. MTE-12 argues that moral disagreement is no deeper than disagreement in other theoretical domains that are not thought to lack facts, and much of it is traceable to disagreement about non-moral facts or to partiality, which are departures from the idealized conditions.
\end{reply}

\begin{objection}[Debunking]
Moral intuitions were shaped by natural selection, which favored reactions that promoted reproductive success, not reactions that track moral truth. Their reliability would therefore be a coincidence \citep{Street2006}.
\end{objection}

\begin{reply}
MTE-09 argues that evolution selected for tracking reasons, and that those reasons include moral ones. The present account adds that the moral facts are idealized reactions to the very features to which evolved reactions respond, so that there is no coincidence to explain: the faculty and its subject matter are related as actual and idealized reactions to the same features. Where evolution has introduced distortions, such as the restriction of concern to the in-group, these are departures from equal attention, and Corollary~\ref{c:fall} classifies the resulting intuitions as unreliable.
\end{reply}

\begin{objection}[Causal idleness]
Perception is reliable because the objects perceived cause the perceptions. Moral facts cause nothing, so the comparison with perception fails \citep{Harman1977}.
\end{objection}

\begin{reply}
On Assumption~\ref{a:ideal}, moral facts are natural facts, and the natural features that constitute them, the suffering of the cat, the cruelty of the boys, are among the causes of the observer's intuition. \citet{Sturgeon1985} argued that moral facts figure in good explanations, as when the injustice of slavery is cited to explain the rise of opposition to it. The objection assumes that moral facts are not natural facts, and our account denies this.
\end{reply}

\section{What the result does not establish}

Theorem~\ref{t:main} depends on Assumption~\ref{a:ideal}, which is defended in MTE-01, and on Assumption~\ref{a:cont}, which is in part empirical. The result does not say how good the conditions must be for a given case, or provide a test by which a subject can tell, from the inside, whether her conditions are good enough; it says only that, when they are, a true intuition yields knowledge. It does not establish that any particular moral principle is known, and it does not establish that intuition is the only source of moral knowledge: inference from known moral premises, and testimony from those who know, extend it. Finally, it does not claim that moral intuitions are perceptual experiences in the strict sense.

\section{Conclusion}

We know moral truths through intuition, a faculty of reaction to particular cases. Moral intuitions are seemings, and so they justify; moral facts are the idealized counterparts of the reactions in which they consist, and so, in good enough conditions, they are safe and yield knowledge; in poor conditions, through ignorance, partiality or framing, they justify but may fail to yield knowledge. The faculty is fallible but generally reliable, and its epistemology is that of perception under imperfect conditions. The knowledge it gives is a posteriori, which is why it does not conflict with the restricted scope of the a priori.

\begin{thebibliography}{99}\small
\bibitem[Audi(2004)]{Audi2004} Audi, R. (2004). \emph{The Good in the Right: A Theory of Intuition and Intrinsic Value}. Princeton University Press.
\bibitem[Audi(2013)]{Audi2013} Audi, R. (2013). \emph{Moral Perception}. Princeton University Press.
\bibitem[Cullison(2010)]{Cullison2010} Cullison, A. (2010). Moral perception. \emph{European Journal of Philosophy}, 18(2), 159--175.
\bibitem[Daniels(1979)]{Daniels1979} Daniels, N. (1979). Wide reflective equilibrium and theory acceptance in ethics. \emph{The Journal of Philosophy}, 76(5), 256--282.
\bibitem[Haidt(2001)]{Haidt2001} Haidt, J. (2001). The emotional dog and its rational tail: A social intuitionist approach to moral judgment. \emph{Psychological Review}, 108(4), 814--834.
\bibitem[Harman(1977)]{Harman1977} Harman, G. (1977). \emph{The Nature of Morality}. Oxford University Press.
\bibitem[Huemer(2005)]{Huemer2005} Huemer, M. (2005). \emph{Ethical Intuitionism}. Palgrave Macmillan.
\bibitem[McBrayer(2010)]{McBrayer2010} McBrayer, J.~P. (2010). A limited defense of moral perception. \emph{Philosophical Studies}, 149(3), 305--320.
\bibitem[McGrath(2004)]{McGrath2004} McGrath, S. (2004). Moral knowledge by perception. \emph{Philosophical Perspectives}, 18, 209--228.
\bibitem[Petrinovich and O'Neill(1996)]{Petrinovich1996} Petrinovich, L. and O'Neill, P. (1996). Influence of wording and framing effects on moral intuitions. \emph{Ethology and Sociobiology}, 17(3), 145--171.
\bibitem[Rawls(1951)]{Rawls1951} Rawls, J. (1951). Outline of a decision procedure for ethics. \emph{The Philosophical Review}, 60(2), 177--197.
\bibitem[Ross(1930)]{Ross1930} Ross, W.~D. (1930). \emph{The Right and the Good}. Oxford University Press.
\bibitem[Sidgwick(1907)]{Sidgwick1907} Sidgwick, H. (1907). \emph{The Methods of Ethics} (7th ed.). Macmillan.
\bibitem[Sinnott-Armstrong(2008)]{SinnottArmstrong2008} Sinnott-Armstrong, W. (2008). Framing moral intuitions. In W.~Sinnott-Armstrong (Ed.), \emph{Moral Psychology}, vol.~2: \emph{The Cognitive Science of Morality}. MIT Press.
\bibitem[Street(2006)]{Street2006} Street, S. (2006). A Darwinian dilemma for realist theories of value. \emph{Philosophical Studies}, 127(1), 109--166.
\bibitem[Sturgeon(1985)]{Sturgeon1985} Sturgeon, N.~L. (1985). Moral explanations. In D.~Copp and D.~Zimmerman (Eds.), \emph{Morality, Reason and Truth}. Rowman \& Allanheld.
\bibitem[V{\"a}yrynen(2018)]{Vayrynen2018} V{\"a}yrynen, P. (2018). Doubts about moral perception. In A.~Bergqvist and R.~Cowan (Eds.), \emph{Evaluative Perception}. Oxford University Press.
\end{thebibliography}

\end{document}
